\section{Homework 13}

\subsection{系统的频率特性}

\subsubsection{根据系统函数绘制系统频率特性}

如图\ref{fig:13-1-1-1}和图\ref{fig:13-1-1-2}所示.

\begin{figure*}[htbp]
    \centering
    \includegraphics[width=0.5\textwidth]{images/13-1-1-1.png}
    \caption{13-1-1(1)(2)(3)波形}\label{fig:13-1-1-1}
\end{figure*}

\begin{figure*}[htbp]
    \centering
    \includegraphics[width=0.5\textwidth]{images/13-1-1-2.png}
    \caption{13-1-1(4)(5)(6)波形}\label{fig:13-1-1-2}
\end{figure*}

\subsubsection{根据系统的零极点分布确定系统的频率特性}

如图\ref{fig:13-1-2-1}, 图\ref{fig:13-1-2-2}, 图\ref{fig:13-1-2-3}, 图\ref{fig:13-1-2-4} 所示.

\begin{figure*}[htbp]
    \centering
    \includegraphics[width=0.5\textwidth]{images/13-1-2-1.png}
    \caption{13-1-2(1)(2)波形}\label{fig:13-1-2-1}
\end{figure*}

\begin{figure*}[htbp]
    \centering
    \includegraphics[width=0.5\textwidth]{images/13-1-2-2.png}
    \caption{13-1-2(3)(4)波形}\label{fig:13-1-2-2}
\end{figure*}

\begin{figure*}[htbp]
    \centering
    \includegraphics[width=0.5\textwidth]{images/13-1-2-3.png}
    \caption{13-1-2(5)(6)波形}\label{fig:13-1-2-3}
\end{figure*}

\begin{figure*}[htbp]
    \centering
    \includegraphics[width=0.5\textwidth]{images/13-1-2-4.png}
    \caption{13-1-2(7)(8)波形}\label{fig:13-1-2-4}
\end{figure*}

\subsubsection{根据零极点分布判断滤波器幅频特性}

(1) 带通滤波器

(2) 高通滤波器

(3) 带通+带阻滤波器

(4) 全通滤波器

(5) 带通滤波器

(6) 带通滤波器

(7) 低通滤波器

(8) 带通滤波器

\subsection{判断最小相位系统}

(1) 不是最小相位系统, 因为有一个零点在右侧.

(2) 不是最小相位系统, 因为在开始时存在过冲

(3) 是最小相位系统, 因为相频特性非单调.

(4) 是最小相位系统, 因为零极点均在单位圆内.

